% Author: Cristian Gonzales
%
% Content only. Layout lives in resume-style.sty.
%
% Editing surface:
%   \cvname{Name}{Subtitle}
%   \cvcontacts{item}{item}{item}{item}
%   \cvsection{Heading}
%   \begin{cvjob}{Company}{Role}{Dates} ... \cvitem{...} ... \end{cvjob}
%   \begin{cvjobdual}{Company}{Title}{Dates}{Prior title}{Prior dates} ... \end{cvjobdual}
%   \cveducation{Institution}{Degree}{Dates}
%   \cvskill{Category}{items}
%   \cvfooter{text}

\documentclass{res}
\usepackage{resume-style}

\begin{document}
\begin{resume}

% ---------------------------------------------------------------------------
% Header
% ---------------------------------------------------------------------------

\cvname{Cristian Gonzales}{Senior Engineer at Microsoft}

\cvcontacts
  {\cvcontact{\faLinkedinSquare}{https://bit.ly/eey7W}{https://bit.ly/eey7W}}
  {\cvcontact{\faGithub}{https://bit.ly/2McyUX}{https://bit.ly/2McyUX}}
  {\cvcontact{\faGlobe}{https://bit.ly/2vA649}{https://bit.ly/2vA649}}
  {\cvcontact{\faEnvelope}{mailto:xcristian.gonzales@gmail.com}{xcristian.gonzales@gmail.com}}

% ---------------------------------------------------------------------------
% Experience
% ---------------------------------------------------------------------------

\cvsection{Experience}

\begin{cvjobdual}{Microsoft Corporation}
  {Senior Software Engineer, Office Data Flywheel}{September 2026 -- Present}
  {Software Engineer, Excel AI}{April 2025 -- September 2026}
  \cvitem{Sole, original author, and owner of the Office
    \cvlink{https://www.nvidia.com/en-us/glossary/data-flywheel/}{Data Flywheel}
    dashboard, a self-service Power BI reporting layer built using
    \cvlink{https://learn.microsoft.com/en-us/kusto/query/}{an internal SQL-like analytics language}
    that unifies Word, Excel, PowerPoint, and Cowork Agent product telemetry
    derived from massive internal datasets processed through
    \cvlink{https://learn.microsoft.com/en-us/azure/machine-learning/concept-ml-pipelines}{Azure Machine Learning (AML) pipelines}.}
  \cvitem{Built \cvlink{https://plotly.com/python/}{Plotly}-driven visualizations
    for user intents, document taxonomies, model comparisons, and OfficeJS
    code-generation diagnostics within the agents, replacing ad hoc query
    exploration with a prioritized product fix backlog and direct input to
    model-selection and vendor decisions.}
  \cvitem{Drove internal adoption to tens of DAU/WAU, and hundreds of distinct users
    (accounting for 90\% of the cluster's queries), streamlining and compressing
    analysis from days to minutes; demoed to Satya Nadella (CEO) through PM
    leadership and set the precedent for Data Flywheel efforts across the
    business unit, with CVP-level visibility.}
  \cvitem{Rated 200/200 for FY26 performance (highest possible), placing in the
    top 1--5\% of engineers at my level across the company.}
  \cvitem{Maintained the
    \cvlink{https://www.microsoft.com/en-us/research/group/prose/}{PROSE research group}'s
    evaluation framework for Excel Agent Mode and built
    Python/\cvlink{https://github.com/microsoft/cpp_client_telemetry}{C++~1DS}
    telemetry bindings to plumb Excel Agent evaluation telemetry into internal
    databases for the purpose of streamlining ingestion and building
    visualizations, giving researchers a consistent view of agent configuration
    changes and ship decisions.}
\end{cvjobdual}

\begin{cvjob}{Microsoft Corporation}{Software Engineer, Office Copilot (Backend)}{December 2024 -- April 2025}
  \cvitem{Prototyped an AI-enriched service to serve dynamic, personalized prompts
    that live within the
    \cvlink{https://support.microsoft.com/en-us/office/chat-with-copilot-about-your-word-document-4482c688-a495-4571-bfcd-4a9fc6608090}{Copilot chat pane}
    of the Office web document editor, to drive habitual usage of Copilot within
    the Office web apps, whereby the LLM prompts are contextualized with
    information about the user from Microsoft Graph and internal Office data
    sources.}
  \cvitem{Core contributor to an experiment, in an effort to further drive
    habituated Copilot usage, to extend the functionality of the
    \cvlink{https://support.microsoft.com/en-us/office/elevate-your-content-with-copilot-in-word-923d9763-f896-4da7-8a3f-5b12c3bfc475}{AI rewrite}
    feature within the Word web editor, to better inform the AI-suggested
    rewrites using previously authored text samples from the user.}
\end{cvjob}

\begin{cvjob}{Microsoft Corporation}{Software Engineer, Office Online (Backend)}{May 2023 -- December 2024}
  % This item uses \item rather than \cvitem: one URL below contains a
  % literal '#', and a '#' inside a macro argument cannot reach \href.
  \item Changed
    \cvlink{https://onenote.officeapps.live.com/hosting/discovery}{Discovery},
    an integral part of the
    \cvlink{https://en.wikipedia.org/wiki/Web_Application_Open_Platform_Interface}{WOPI}
    protocol that powers \textit{all} of Office Online, to pare down
    \cvlink{https://learn.microsoft.com/en-us/microsoft-365/cloud-storage-partner-program/online/discovery}{Office Online's partner storage integration offering}
    to only the
    % \href is used directly here because the URL contains a literal '#',
    % which cannot pass through \cvlink's argument.
    \href{https://learn.microsoft.com/en-us/microsoft-365/cloud-storage-partner-program/online/discovery#wopi-actions}{officially supported Office apps and app capabilities for partners},
    along with mapping the partner test version of Discovery to a more
    production-ready version--this yielded the net result of not exposing
    internal prototyping and test datacenters, which can prove to be extremely
    volatile, to real customers (i.e. an outage internally impacted external
    customers).
  \cvitem{As part of shipping new, CDN-deployed customer funnels to the flagship
    Office apps, found at
    \cvlink{https://powerpoint.cloud.microsoft}{powerpoint.},
    \cvlink{https://excel.cloud.microsoft}{excel.}, and
    \cvlink{https://word.cloud.microsoft}{word.cloud.microsoft}, leveraged
    existing Office Online infrastructure to address a feature gap in an
    internal offering of Azure Front Door (AFD), a layer-7 aware traffic
    manager, by building a service on the Office Online stack to absorb failed
    requests from AFD and proxy them to the CDN to serve the static content in
    question, improving the reliability of serving the CDN-deployed content by
    2\%.}
  \cvitem{As a byproduct of standing up new web services within Office Online's
    infrastructure to support
    \cvlink{https://techcommunity.microsoft.com/t5/microsoft-365-blog/introducing-cloud-microsoft-a-unified-domain-for-microsoft-365/ba-p/3804961}{Office's effort to bring Office to the cloud.microsoft domain},
    introduced the usage of the .NET 6+ runtime and Kestrel as a web server, in
    lieu of .NET Framework and Internet Information Services (IIS), for which
    there was precedent for, which yielded a ballpark savings of 20\% in
    compute costs.}
  \cvitem{Migrated server-to-server authentication flows for critical internal
    proxy services supporting Office Online, from using app-asserted Bearer
    tokens to more secure Proof-of-Possession tokens, mitigating the risk of
    high-privilege access (Bearer) token theft and unauthorized access to
    protected resources, that of which was a top-down C-suite driven effort to
    enhance customer trust across Microsoft.}
\end{cvjob}

\begin{cvjob}{Microsoft Corporation}{Software Engineer, Microsoft Loop (Backend)}{May 2022 -- May 2023}
  \cvitem{Led the design and implementation of a waitlist management API, using a
    proprietary NoSQL persistent store, for
    \cvlink{https://loop.cloud.microsoft/learn}{Loop}, a new product in the
    Office 365 ecosystem, with an upper bound of 200,000 new users at public
    preview.}
  \cvitem{Implemented European Union data boundary (EUDB) -compliant merging
    logic for Loop user ``presences'' in cross-boundary ``workspaces'', and
    EUDB-compliant routing of requests to datacenters nearest to the client,
    bringing down the P99 latency times for a single user presence request from
    3 seconds to a few hundred milliseconds.}
  \cvitem{Implemented server-to-server (S2S) authentication, using Azure Active
    Directory (AAD) tokens, for Loop and reduced latencies using on-box token
    caching, bringing down the P99 latency for some scenarios by more than
    50\%.}
\end{cvjob}

\begin{cvjob}{Social Finance, Inc. (SoFi)}{Senior Software Engineer (Backend)}{March 2022 -- May 2022}
  \cvitem{Contributor to a service adapter that served as a read-only caching
    layer to all microservices across the company so they could meet their
    SLAs, in an effort to reduce P99 read times for some services in the order
    of hundreds of milliseconds.}
  \cvitem{Prototyped a Debezium listener that listened to data change events from
    event sources onboarded to our service, for the purpose of dynamically
    storing them in DynamoDB, regardless of the shape (schema) of the Debezium
    event payload.}
  \cvitem{Introduced code infrastructure to support LocalStack in development
    workflows to avoid the overhead of redeploying AWS infrastructure to our
    staging environment, allowing for leaner workflows where redeployments took
    seconds instead of minutes.}
\end{cvjob}

\begin{cvjob}{Meredith Corporation (Bizrate Insights)}{Software Engineer (Backend)}{July 2021 -- March 2022}
  \cvitem{Introduced in-memory caches to microservices (AWS Lambda) to reduce
    network requests to an internal A/B testing service, reducing response
    times of microservices of interest by roughly 80\% for ``hot'' Lambda
    containers--benchmarks in the order of 5 seconds to a few hundred
    milliseconds.}
  \cvitem{Minified, and introduced CDNs (AWS CloudFront) for, static content
    being served by a core Java Spring microservice, cutting down the time to
    serve static content by \textgreater 80\% in some cases.}
  \cvitem{Architected a management microservice (AWS Lambda) sitting behind a
    REST controller (AWS API Gateway) to ban users at the network boundary,
    using AWS Cloud Development Kit and AWS Web Application Firewall (WAF),
    whereby WAF would examine universally unique identifiers associated with the
    user.}
  \cvitem{Introduced Object Relational Mapping, using SQLAlchemy, in the form of
    a Data Access Layer and packaged as a Python pip distribution, to replace ad
    hoc raw database operations being made in Apache Airflow jobs.}
  \cvitem{Cultivated a strong engineering culture by championing Test-Driven
    Development where there was otherwise none, resulting in the turnaround time
    for features in the Software Development Lifecycle (SDLC) being cut by
    nearly half, as bugs where caught earlier in the SDLC, and SDETs and
    developers were able to work in parallel rather than serially.}
  \cvitem{Championed DevOps by defining cloud infrastructure-as-code through
    CloudFormation, allowing for infrastructure to become less anecdotal and
    helping developer velocity when an engineer would try to understand the
    surrounding infrastructure of a given part of the greater architecture.}
\end{cvjob}

\begin{cvjob}{Northrop Grumman}{Software Engineer (Infrastructure)}{August 2020 -- July 2021}
  \cvitem{Developed a library of dynamic HashiCorp Packer configurations to build
    DISA and NIST STIG-hardened RHEL and CentOS machine images for VMWare
    VSphere and Docker, using Gradle (Groovy) and an in-house Ansible library,
    migrating deployments from PXE-boot to template-based and reducing
    deployment times by roughly 60\%.}
  \cvitem{Developed Jenkins pipelines to deploy hardened infrastructure and test
    in-house Ansible configurations using Docker workflows, allowing the team to
    save at least 80 developer hours per month, and developed VSphere wrapper
    APIs, using Java and Gradle, for Jenkins resource management.}
  \cvitem{Wrote Ansible playbooks to harden firewalld with respect to NIST STIGs,
    and tested them by spoofing IP-layer header attributes and creating Martian
    packets, using Python's Berkeley Sockets API interface and kernel tunables.}
  \cvitem{Wrote Ansible roles to configure and harden deployments of Cisco NX-OS
    Layer-3 switches, VMWare ESXi guests, Linux NetworkManager, and Chrony
    timeservers.}
\end{cvjob}

\begin{cvjob}{Northrop Grumman}{Software Engineer (Build and Release)}{July 2019 -- August 2020}
  \cvitem{Developed Jenkins pipelines to streamline different build processes in
    Linux, Solaris, and Windows, decreasing the time of the build and release
    process from four days to two hours,
    \cvlink{https://web.archive.org/web/20230204223645/https://news.northropgrumman.com/news/features/northrop-grumman-increases-collaboration-with-space-and-missile-systems-center-by-implementing-agile-methodology}{recognized in a company press release}.}
  \cvitem{Ported static Jenkins web interface logic into a dedicated Groovy DSL
    codebase to describe job configurations as code, tailored for the needs of
    different Jenkins pipeline workflows (developer sandboxes, release builds,
    nightly builds, etc.), where efforts allowed for an estimated 40\% decrease
    of COGS for the configuration management team.}
  \cvitem{Automated deployment for builds for ad hoc projects and ported build
    systems from Apache Ant to Apache Maven.}
\end{cvjob}

\begin{cvjob}{Northrop Grumman}{Software Engineer (Backend)}{July 2019 -- August 2020}
  \cvitem{Developed an API to filter JMS messages from a legacy JMS service,
    using Guice injections to differentiate between different instances of the
    service and JAXB for object serialization to send messages across the wire.}
  \cvitem{Refactored database schemas and their associated JPA mappings, along
    with protocol buffer definitions, to differentiate between different
    satellite missions for a server-side event message syslog implementation,
    where the mission identifier was injected using Guice.}
  \cvitem{Developed a JAXB unmarshalling utility so subscribers may filter syslog
    event messages based on the satellite mission, where subscriptions were made
    through the unmarshalled data and Guice dependency injections.}
  \cvitem{Served as a Scrum Master, leading daily Scrum meetings, retrospectives,
    and planning.}
\end{cvjob}

\begin{cvjob}{Northrop Grumman}{Software Engineering Intern}{June 2017 -- August 2017}
  % \cvitem{Interfaced with NASA's GMSEC API to build a visualization tool which tracked health statuses of satellites.}
  % \cvitem{Ported payload data, via web sockets, across multiple Northrop Grumman domains/applications while leveraging code reuse through dependency injections, using Google Guice.}
  \cvitem{Prototyped a visualization tool to track satellite health statuses with
    NASA's GMSEC API and Guice injections.}
\end{cvjob}

% ---------------------------------------------------------------------------
% Education
% ---------------------------------------------------------------------------

\cvsection{Education}

\cveducation{University of California, Santa Cruz}{Computer Science, B.S.}{June 2016 -- June 2019}

\cveducation[Selected through leadership nomination to be part of Microsoft Office Product Group's (OPG) Applied Science Upskilling Program for Engineers, with a capstone project at the conclusion of the course.]{Stanford School of Engineering: \cvlink{https://online.stanford.edu/programs/artificial-intelligence-professional-program}{Artificial Intelligence Professional Program}}{\cvlink{https://online.stanford.edu/courses/xcs221-artificial-intelligence-principles-and-techniques}{XCS 221: Artificial Intelligence: Principles and Techniques}}{August 2026 -- December 2026}

% ---------------------------------------------------------------------------
% Skills
% ---------------------------------------------------------------------------

\cvsection{Skills}

\begin{cvskills}
  \cvskill{Languages}{Bash, C, C\#, Go, Java, Python, SQL, TypeScript}
  \cvskill{Tools and ecosystems}{Apache Airflow, AWS, Docker, Git, Hugging Face, \LaTeX, PyTorch}
  \cvskill{Concepts}{Applied AI, Build and Release Engineering, Computer Networking, Containerization and Virtualization, Deep Learning, DevOps, Distributed Systems, Experimentation and A/B Testing, Machine Learning, Operating Systems}
\end{cvskills}

\end{resume}

\cvfooter{References available upon request.}

\end{document}
